\chapter[Préparer une campagne d’annotation]{Préparer une campagne d’annotation experte : la méthodologie de la fabrique des métadonnées}

\lettrine{L}{'une des} missions de ce stage a été d'entraîner un modèle généraliste. Par \og généraliste \fg, nous entendons un modèle capable d'extraire des éléments visuels indifféremment au sein de documents manuscrits ou imprimés. Ce défi implique de surmonter la forte hétérogénéité visuelle de ces supports, allant de la mise en page normée des débuts de l'imprimerie aux structures plus fluides et parfois altérées des manuscrits médiévaux. L'objectif est d'aboutir à une ontologie suffisamment unifiée et opérationnelle pour s'affranchir du bruit visuel comme des spécificités propres à une époque donnée.

\section[L'ontologie en machine learning]{Qu'est-ce qu'une ontologie en machine learning ?}

\subsection{L'importance de l'ontologie pour les méthodes d'apprentissage automatique}

En vision artificielle, une ontologie d'annotation désigne le cadre conceptuel et structuré qui définit la manière dont les données brutes (ici, des images de documents historiques) doivent être catégorisées et étiquetées pour entraîner un modèle d'apprentissage automatique. Loin de se réduire à une simple liste de mots-clés appliqués sur des boîtes englobantes, il s'agit d'un objet complexe qui s'articule autour de trois dimensions fondamentales :

\begin{itemize}
	\item \textbf{Le vocabulaire (ou taxonomie) :} l'ensemble fini et contrôlé des classes retenues. Dans notre cadre, il réunit les classes \textit{Illustration}, \textit{Lettrine}, \textit{Tampon}, \textit{Ornementation} et \textit{Table}.
	
	\item \textbf{La sémantique :} l'ensemble des règles d'annotation qui définissent la tâche d'annotation, le périmètre d'application du vocabulaire et ses exceptions.
	
	\item \textbf{La structure :} l'organisation des relations entre les classes. Une ontologie peut être hiérarchique (comme le prévoient des standards tels que SegmOnto) ou plate (sans hiérarchie).
\end{itemize}

Le but de la création d'une ontologie est de traduire la complexité et la réalité d'un document historique en un langage fini, compréhensible et traitable tant par des annotateurs humains que par un algorithme. Lors de notre stage, l'enjeu autour de la définition de cette ontologie était double : d'une part, elle influençait directement la manière dont notre modèle allait être réintégré et exploité par les équipes d'AIKON sur l'application ; d'autre part, l'ontologie sous-jacente à \textit{GenHisDoc} constituait l'élément conditionnant la réutilisation de certains jeux de données au détriment d'autres, en fonction de leur schéma d'origine et de la forme technique de leurs annotations. De plus, l'emploi d'une ontologie trop complexe, qui nécessiterait une infrastructure technique lourde, réduirait la facilité de réutilisation des données produites.

\subsection{Les ontologies appliquées aux documents historiques}

Comme nous l'avons souligné précédemment, la plateforme AIKON exploite encore des modèles de détection monoclasses. Une part importante de notre travail a donc consisté à tester la mise en place d'un modèle à la fois généraliste, capable de traiter une grande variété de types de documents, et multiclasse. Dès lors, nous n'avons pas hérité d'une ontologie rigidement fixée plus tôt dans le projet.

\subsubsection{Les ontologies PAGE XML et ALTO XML}

En tant que schémas XML, \gls{page} et \gls{alto} embarquent une ontologie structurelle de la page et constituent aujourd'hui les standards d'export privilégiés par les institutions patrimoniales. Ils sont notamment employés dans le cadre des jeux de données de référence comme ceux du laboratoire \textit{PRImA}, utilisés lors des compétitions internationales de reconnaissance de documents telles que l'\gls{icpr} ou \gls{icdar}.

Dans le détail, \gls{page}\footnote{S. Pletschacher, A. Antonacopoulos, \og The PAGE (Page Analysis and Ground-Truth Elements) Format Framework \fg, \textit{20th International Conference on Pattern Recognition (ICPR2010)}, IEEE-CS Press, pp. 257-260, 2010, DOI : \href{https://doi.org/10.1109/ICPR.2010.72}{10.1109/ICPR.2010.72}.} (\textit{Page Analysis and Ground-Truth Elements}) est développé par le \textit{PRImA research lab} de l'université de Salford (Manchester) et repose sur une architecture hiérarchique. Son ontologie distingue d'abord de grandes régions de blocs (\textit{TextRegion}, \textit{ImageRegion}, \textit{MathsRegion}, \textit{SeparatorRegion}), qui peuvent elles-mêmes être subdivisées en lignes de texte, en mots, puis en glyphes individuels. \gls{page} associe à ses classes des coordonnées géométriques complexes pour des tâches de segmentation fine, ainsi que des métadonnées internes (langue, orientation, typographie).

De son côté, \gls{alto} est un standard maintenu à l'origine grâce au financement européen du projet METAe\footnote{B. Stehno, A. Egger, G. Retti, \og METAe --- Automated Encoding of Digitized Texts \fg, \textit{Literary and Linguistic Computing}, vol. 18, n° 1, pp. 77–88, 2003, DOI : \href{https://doi.org/10.1093/llc/18.1.77}{10.1093/llc/18.1.77}.}. Souvent associé au format de métadonnées de conservation \gls{mets}, il constitue le standard privilégié par de grandes institutions patrimoniales internationales telles que la \gls{bnf} (Gallica) ou la \textit{Library of Congress} qui héberge le schéma officiel depuis 2009. Il se concentre sur la description physique et logique de la page, structurant les blocs de texte, les illustrations et les espaces blancs grâce à des coordonnées exprimées sous forme de boîtes englobantes.

Cependant, si la richesse sémantique et structurelle de ces deux standards garantit une description exhaustive des documents historiques, elle présente une limite majeure dans le cadre technique de notre travail. Ces ontologies sont profondément imbriquées et hiérarchiques. Or, pour entraîner des architectures de détection d'objets rapides et efficaces comme les modèles \gls{yolo} intégrés à AIKON, les données d'entrée doivent idéalement reposer sur une ontologie \og plate \fg. Si certains projets d'envergure disposent des ressources nécessaires pour maintenir des structures aussi complexes, nous avons fait le choix de nous en écarter afin de privilégier un modèle d'annotation direct et aisément convertible.


\subsubsection{SegmOnto}

Au début de notre réflexion, nous nous sommes appuyés sur SegmOnto\footnote{S. Gabay, A. Pinche, K. Christensen, J.-B. Camps, \og SegmOnto: A Controlled Vocabulary to Describe and Process Digital Facsimiles \fg, \textit{Journal of Data Mining and Digital Humanities}, 2024, DOI : \href{https://doi.org/10.46298/jdmdh.12689}{10.46298/jdmdh.12689}.}, l'ontologie de référence pour la description des documents historiques. Si nous en avons conservé la philosophie globale, ancrée dans la description morphologique de la page, nous avons dû nous en écarter pour répondre à une double contrainte, à la fois technique et scientifique.

Sur le plan technique, l'architecture des modèles \gls{yolo} nous imposait de renoncer à toute arborescence complexe au profit d'une ontologie strictement « plate ». Par ailleurs, pour garantir la robustesse de l'apprentissage et éviter un éparpillement statistique des détections, nous avons fait le choix de réduire le schéma à un maximum de cinq classes essentielles.

Sur le plan scientifique, les objectifs du projet AIKON exigeaient un degré de granularité qui entrait en contradiction avec certaines agrégations de SegmOnto. C'est notamment le cas de la classe \textit{GraphicZone}, qui fusionne toutes les expressions picturales. Pour répondre aux besoins de dépouillement historique, nous avons dû réintroduire une distinction fonctionnelle en scindant cette catégorie en deux classes distinctes (\textit{Illustration} et \textit{Ornementation}), en nous inspirant de l'ontologie du jeu de données S-VED.

\begin{table}[h]
    \centering
    \ra{1.2}
    \caption{L'ontologie SegmOnto}
    \label{tab:segmonto}
    \begin{tabular}{l l}
        \toprule           
        \textbf{Classes SegmOnto} & \textbf{Définition} \\
        \midrule
        CustomZone & ... \\
        DamageZone & Dégâts sur le document \\
        DigitisationArtefactZone & Dégâts sur l'image numérisée \\
        DropCapitalZone & Lettrine \\
        GraphicZone & Dessins, Illustration, Ornementation \\
        MainZone & Texte principal au niveau de la page \\
        MarginTextZone & Textes marginaux \\
        MusicZone & Partitions \\
        NumberingZone & Pagination \\
        QuireMarksZone & Réclame \\
        RunningTitleZone & Titres courants \\
        SealZone & Sceaux \\
        StampZone & Estampilles \\
        TableZone & Tables de données \\
        TitlePageZone & Page de titre \\
        \bottomrule
    \end{tabular}
\end{table}

\section[Construire une ontologie d'annotation]{Construire une ontologie d'annotation : de la théorie à la pratique}

Une ontologie aussi synthétique que la nôtre (limitée à cinq classes) déplace la complexité du modèle de représentation des données vers la phase d'annotation manuelle. Il est donc devenu indispensable d'établir des consignes d'annotation strictes afin de départager des éléments graphiquement proches mais fonctionnellement distincts au sein des documents historiques.

\subsection{Définir les frontières sémantiques : le défi des cas limites}

La principale difficulté a résidé dans la distinction entre les éléments relevant de la classe \textit{Illustration} et ceux situés à la frontière entre \textit{Illustration} et \textit{Ornementation}. Dans les manuscrits médiévaux comme dans les premiers imprimés, la frontière entre le décoratif et le figuratif est souvent poreuse. Nous avons fait le choix de nous appuyer sur le rapport au texte : une \textit{Illustration} possède une fonction sémantique ou didactique directe (elle illustre explicitement le propos), tandis que l'\textit{Ornementation} (tels que les bandeaux, culs-de-lampe ou frises marginales) remplit une fonction purement esthétique ou de séparation structurelle.

De la même manière, la classe \textit{Initial} a exigé des règles de délimitation claires. Une lettrine historiée (contenant une scène figurative) devait-elle être annotée comme une illustration ou comme une lettrine ? Pour maintenir la cohérence de l'apprentissage du modèle, nous avons décidé de prioriser la fonction typographique : tout caractère alphabétique agrandi marquant le début d'une section est étiqueté \textit{Initial}, indépendamment de sa richesse iconographique. La prise en compte de cette classe a également nécessité l'étude de plusieurs cas hybrides, notamment des lettres isolées ou d'autres signes typographiques comme le pied-de-mouche (¶), susceptibles d'être confondus avec le motif cible.

La définition de la classe \textit{Table} a elle aussi été source de problèmes. Contrairement aux documents contemporains où les données tabulaires sont généralement structurées par des grilles ou des lignes séparatrices explicites, les documents historiques s'appuient fréquemment sur des structures implicites (alignements par des espaces, listes structurées). Il a fallu déterminer à partir de quel degré de structuration spatiale un bloc de texte devient une \textit{Table}. Nous avons pris le parti de caractériser cette classe par la présence évidente d'une relation bidimensionnelle (lignes et colonnes) porteuse de sens, même en l'absence de lignes de séparation tracées. Pour aider le modèle à repérer cette régularité spatiale, la consigne a été donnée d'englober l'intégralité de la structure logique du tableau, incluant ses en-têtes de colonnes lorsqu'ils sont présents. Enfin, nous avons fait le choix d'exclure de cette classe les structures d'organisation globale du document, telles que les tables des matières ou les index.

\begin{table}[h]
    \centering
    \ra{1.2}
    \caption{L'ontologie de GenHisDoc}
    \label{tab:segmonto}
    \begin{tabular}{l l l}
        \toprule           
        \textbf{Label Yolo} & \textbf{Classes GenHisDoc} & \textbf{Définition} \\
        \midrule
        0 & Illustration & Visuel qui se réfère au texte \\
        1 & Initial & Lettrine \\
        2 & Ornament &  Visuel qui ne se réfère pas au texte \\
        3 & Stamp & Estampille de bibliothèque \\ 
        4 & Table & Données tabulaires \\
        \bottomrule
    \end{tabular}
\end{table}

\subsection{Mettre en place l'ontologie : écrire un guide d'annotation}

Comme le montre l'organigramme~\ref{fig:GenHisDoc}, la rédaction du guide d'annotation a été réalisée \textit{a posteriori} du travail d'annotation et de correction initial. GenHisDoc n'étant conçu que comme un test préalable à une campagne d'annotation de plus grande envergure, la construction de l'ontologie a suivi un processus empirique et itératif : nous avons d'abord annoté un premier volume de données sur divers types de documents, puis arrêté des choix clairs quant à la structure de l'ontologie, avant de procéder à une phase de rétro-correction afin d'harmoniser les données antérieures avant l'entraînement des modèles.

Le but d'un guide d'annotation est de fournir au modèle des données d'apprentissage parfaitement cohérentes grâce à un fort accord inter-annotateurs, tant sur le plan géométrique (précision des délimitations) que sur le respect des règles et exceptions sémantiques propres à chaque classe au sein du document. Dans le cadre d'AIKON, ce guide constitue également un document méthodologique essentiel permettant de conserver une trace explicite des principes retenus, des retours d'expérience et des divers cas limites rencontrés lors de la constitution du jeu de données.

Le guide d'annotation s'articule autour de deux volets principaux. Le premier présente le contexte du projet, définit la tâche d'annotation attribuée et pose les règles générales applicables à l'ensemble des opérations de délimitation.


\subsubsection{Les règles générales}

Les règles générales de cadrage établissent les exigences en matière de rigueur géométrique (c'est-à-dire la méthode précise de tracé de la boîte englobante). Elles imposent notamment un ajustement au plus près des contours de l'objet (\og \textit{hug the border} \fg), la séparation stricte d'éléments visuellement disjoints, ainsi qu'un principe empirique fondamental : l'interdiction d'annoter par déduction ce qui n'est pas directement visible à l'œil humain (\og \textit{annotate what is visible, not what you infer} \fg).

Ces règles générales visent à maximiser l'accord inter-annotateurs et à réduire le bruit lors de l'entraînement des modèles, en évitant d'inclure inutilement de l'arrière-plan ou du texte environnant au sein de la zone délimitée.

\subsubsection{Les règles ontologiques}

D'autre part, la section consacrée aux définitions de classes (\textit{Class definition}) décline, pour chaque étiquette (\textit{Illustration}, \textit{Ornementation}, \textit{Lettrine}, \textit{Tampon}, \textit{Table}), une structure explicite en trois points :

\begin{enumerate}
	\item \textbf{Une description fonctionnelle} rappelant le rôle sémantique et visuel de l'élément au sein de la page ;
	\item \textbf{Des règles d'annotation précises} fixant le périmètre exact de la boîte englobante ;
	\item \textbf{Une analyse des cas limites (\textit{edge cases})} apportant une réponse systématique aux typologies complexes identifiées lors de la phase de test (telles que les lettrines historiées, les encadrements ornés ou les éléments fragmentés).
\end{enumerate}

Outre la rigueur des définitions textuelles, l'usage systématique de figures et la juxtaposition d'exemples concrets (cas validés versus erreurs à éviter) lèvent immédiatement les ambiguïtés. Le format de diffusion de ces consignes s'avère tout aussi déterminant que leur contenu. Pour maintenir un engagement fort et limiter les dérives d'annotation, un guide moderne ne saurait se réduire à un simple document PDF statique : il gagne à être un dispositif dynamique et interactif. Les meilleures directives doivent \og captiver \fg{} l'annotateur dès son intégration grâce à des modules de test ou des quiz d'évaluation, puis s'insérer directement au cœur de l'interface d'annotation afin d'offrir, au simple clic sur une classe, l'affichage instantané de ses cas limites sans jamais quitter l'espace de travail.





